% discussion.tex -- Batch Aggregation Paper (BNR-2026-004)

\section{Discussion}

\subsection{Hypothesis Verdict}

All four components of the hypothesis are confirmed with large margins:

\begin{enumerate}
\item \textbf{Throughput $\geq$ 100 tx/min}: Confirmed. Measured
  14,150--274,438~tx/min across configurations, exceeding the target by
  141$\times$ to 2,744$\times$.
\item \textbf{Batch formation latency $<$ 5,000~ms}: Confirmed. Measured
  0.010--0.022~ms, exceeding the target by over 200,000$\times$.
\item \textbf{Zero transaction loss}: Confirmed empirically (150/150 tests)
  and formally (TLA+ model check PASS after protocol correction).
\item \textbf{Batch determinism}: Confirmed (450/450 tests).
\end{enumerate}

However, the zero-loss claim requires qualification. The \emph{implemented}
protocol (checkpoint at formation) contains a silent loss bug detectable only
through formal verification. The \emph{corrected} protocol (deferred
checkpoint) achieves true zero loss, verified by exhaustive model checking.

\subsection{Trade-off Analysis}

The deferred checkpoint protocol introduces one trade-off: on crash recovery,
the system must re-prove at most one batch of transactions (1.9--12.8 seconds
of redundant computation). This is acceptable given that the mean time between
failures (MTBF) for enterprise server hardware is measured in thousands of
hours, while the re-proving cost is under 13 seconds. The WAL segment also
grows slightly larger (uncommitted entries include batched transactions),
requiring post-checkpoint compaction---already planned as an optimization.

The HYBRID strategy dominates SIZE-only batching (16\% higher throughput,
equal safety) and TIME-only batching (which fails entirely under burst
arrivals). This matches the universal design choice across production
ZK-rollups~\cite{chaliasos2024analyzing}.

\subsection{The Value of Formal Verification}

The central finding of this work is methodological. The TLA+ model checker
discovered a critical safety violation in under one second that 150 empirical
crash recovery tests did not detect. This is not because the test suite was
poorly designed---the 5 crash scenarios represent a reasonable partition of
the failure space. Rather, it demonstrates a fundamental limitation of
example-based testing: tests verify specific execution paths, while model
checking explores \emph{all} interleavings.

The missing scenario (crash between checkpoint and batch processing) is
subtle because it requires reasoning about the interaction between two
design decisions:
\begin{itemize}
\item The WAL checkpoint advances at formation, not processing (a design choice).
\item Formed batches exist only in volatile memory (an implementation fact).
\end{itemize}

Neither fact alone is dangerous. Their \emph{combination} creates a
durability gap. This is precisely the class of bugs that formal methods
excel at finding: emergent properties of component interactions that are
invisible to unit testing.

The severity is high: the loss is \emph{silent} (no error raised), the
window is \emph{long} (seconds, not microseconds), and the consequence
is \emph{irrecoverable} (lost transactions cannot be re-enqueued from the
enterprise source without external intervention). For an enterprise
traceability system where missing records can invalidate audit chains,
this represents a critical safety failure.

\subsection{Limitations}

\begin{enumerate}
\item \textbf{Synchronous benchmark}: All transactions arrive in a tight
  loop (burst mode). Real enterprise workloads have Poisson-like arrival
  patterns. The TIME strategy requires asynchronous evaluation to
  demonstrate its value as a timeout-based fallback.

\item \textbf{Windows fsync semantics}: Windows NTFS may not guarantee
  true \texttt{fdatasync}. Production benchmarks on Linux with ext4/XFS
  and \texttt{O\_DIRECT} will show different WAL write latencies (likely
  higher, closer to the published 880~$\mu$s benchmark).

\item \textbf{Single writer}: The benchmark uses a single writer thread.
  Multi-enterprise scenarios with concurrent writers require partitioned
  queues or lock-free data structures.

\item \textbf{JSON WAL format}: At $>$10K~tx/min, JSON serialization
  overhead (300~bytes per entry vs.\ $\sim$50 bytes binary) becomes
  significant. Production deployments should evaluate binary WAL formats.

\item \textbf{No SMT integration}: The benchmark does not perform Sparse
  Merkle Tree state transitions. In production, batch formation includes
  SMT updates (1.8~ms per insert from RU-V1), adding latency proportional
  to batch size.

\item \textbf{Finite model}: The TLA+ model uses 4 transaction identifiers
  and batch threshold 2. While TLC achieves complete state-space coverage
  for this configuration, the result generalizes by symmetry reduction
  only---larger models may exhibit behaviors not present in the small model.
  The safety properties are structurally independent of the number of
  transactions (they depend on partition invariants, not cardinality).
\end{enumerate}

\subsection{Reconciliation with Literature}

Our WAL write latencies (149--210~$\mu$s) are consistent with published
database WAL benchmarks when accounting for JSON serialization and SHA-256
checksum computation. The 300$\times$ difference from the raw binary
benchmark (529~ns~\cite{chia2024database}) is fully attributable to format
overhead, not implementation inefficiency.

Our throughput (274K~tx/min peak) is consistent with single-process
in-memory queue performance. It falls between RabbitMQ quorum queues
(2K~msg/s with fsync) and Kafka (2M/s distributed). The comparison is
directional---our system targets enterprise validium (single-tenant,
100+~tx/min) rather than public infrastructure.

The HYBRID strategy recommendation aligns with all surveyed production
systems: zkSync Era, Polygon zkEVM, Scroll, and Ethrex all use combined
size-time triggers for batch formation.
